% Generated draft: requires mathematical and source review.
\documentclass[11pt]{article}
\usepackage[margin=25mm]{geometry}
\usepackage{fontspec}
\setmainfont{Latin Modern Roman}
\usepackage{amsmath,amssymb,mathtools,unicode-math}
\setmathfont{Latin Modern Math}
\usepackage{enumitem,needspace,xurl,hyperref}
\setlength{\parindent}{0pt}
\setlength{\parskip}{5pt}
\setlength{\emergencystretch}{3em}
\newcommand{\N}{\mathbb{N}}
\newcommand{\Z}{\mathbb{Z}}
\newcommand{\Q}{\mathbb{Q}}
\newcommand{\R}{\mathbb{R}}
\newcommand{\C}{\mathbb{C}}
\newcommand{\F}{\mathbb{F}}
\newcommand{\A}{\mathbb{A}}
\newcommand{\PP}{\mathbb{P}}
\newcommand{\E}{\mathbb{E}}
\newcommand{\eps}{\varepsilon}
\newcommand{\dd}{\,\mathrm{d}}
\begin{document}
\section*{UnsolvedMath ID: 20001982}
\textbf{Source problem code:} catalogue entry
\subsection{Definitions and mathematical statement}
J. Weitsman's Question 13.1 asks what can be said about quotient cohomology when a compact K\"ahler Hamiltonian manifold has only diagonal Hodge cohomology. In the regular interpretation of the aid, let a compact Lie group \(K\) act holomorphically and Hamiltonianly on \((M,\omega)\), with moment map \(\mu\), and assume \(Z=\mu^{-1}(0)\) is nonempty and the \(K\)-action on it is locally free. Then \(X=Z/K\) is a compact K\"ahler orbifold; a free action gives the smooth case. Use ordinary orbifold differential-form Hodge cohomology, not an unspecified twisted-sector theory, and write
\[
 h^{p,q}(Y)=\dim_{\mathbb C}H^{p,q}(Y).
\]
The diagonal hypothesis is \(h^{p,q}(M)=0\) for \(p\ne q\).

The aid attributes to Yi Lin a Hodge-compatible equivariant Dolbeault Kirwan map, surjective in each bidegree, giving
\[
 h^{p,q}(X)=0\qquad(p\ne q).
\]
It reports this as settling the smooth and locally free orbifold readings, not arbitrary singular quotients. Complete hypotheses and a full citation for that theorem are not printed.

For a torus \(K=T^r\), \(r\geq1\), the aid additionally reports, for every \(k\geq0\),
\[
 h^{k,k}(X)\leq\sum_{i=0}^k h^{i,i}(M)
                  \binom{r+k-i-1}{k-i}.
\]
The degree-two torus generators have Hodge bidegree \((1,1)\), fixing the coefficient convention in this bound. It follows in this setting that odd Betti numbers vanish and
\[
 P_X(t)=\sum_kh^{k,k}(X)t^{2k},\qquad
 H_X(u,v)=\sum_kh^{k,k}(X)(uv)^k.
\]
These are constraints on dimensions; no claim of a full ring presentation or algebraic-cycle generation is made. The rank-zero case is simply the unreduced manifold and is not inserted into the displayed binomial convention.

The source includes two historical comments. R. Sjamaar's Comment 13.2 discusses quantization commuting with reduction for cohomology with values in a prequantum line bundle \(L\), and Teleman's [Tel00] extension involving \(L\otimes\Omega^q\). It suggests that using \(\Omega^q\) alone might force corresponding vanishings downstairs. Such an untwisted conclusion is a proposed route, not justified merely by deleting \(L\) from a theorem with positivity or linearization assumptions. The original writes \(H^{p,q}(M)=H^p(M;\Omega^q)\); with the standard Dolbeault ordering one writes \(H^{p,q}(M)=H^q(M;\Omega^p)\). This index convention does not change the diagonal-vanishing question. Comment 13.3, attributed to D. Burns and J. Weitsman, says they found a method, possibly implicit in earlier work, but supplies no details. Neither comment defines the cohomology theory or conclusion for an arbitrary singular quotient.
\subsection{Short English statement}
For holomorphic Hamiltonian K\"ahler reductions with locally free zero level, the supplied report says diagonal Hodge type descends, via Yi Lin's componentwise-surjective Hodge Kirwan map. It adds an explicit torus Hodge-number bound. Singular quotients and the suggested untwisted quantization argument require further specification.
\subsection{Sources}
\begin{itemize}
\item Local source record: catalogue entry, catalogue entry; canonical \texttt{aim-geometry-notes.json}, \texttt{notes[319]}, source item 13.1. Title: ``Diagonal Hodge type under regular K\"ahler reduction.''
\item AIM workshop list \emph{Moment maps and surjectivity in various geometries}, Question 13.1 attributed to J. Weitsman, Comment 13.2 to R. Sjamaar, and Comment 13.3 to D. Burns and J. Weitsman; supplied reference \url{https://aimath.org/WWN/momentmaps/momentmaps.pdf}. The verification reports comparison with page 16; this draft does not access a PDF.
\item Original reference Teleman [Tel00] for twisted sheaf-cohomology reduction. The supplied summary attributes the Hodge-compatible Kirwan theorem to Yi Lin. Full bibliographic entries and theorem numbers are absent. The coefficientwise torus bound is a candidate contribution of attempt 1.
\item Registry dated 14 August 2026: ``partially\_solved,'' ``PARTIAL-PROGRESS,'' statement status ``exact''; attempt partial, rigor ``proved,'' novelty confidence low. The machine-generated, non-peer-reviewed aid requires independent verification before citation.
\end{itemize}
\end{document}
